% (C) 2026 Ralf Stephan, in collaboration with Claude Code.  CC0 1.0.
\chapter{Labels, languages and sofic shifts}\label{ch:sofic}

Chapter~\ref{ch:entropy} counted labelled paths. Two paths with the same labels were counted
twice. In the applications only the labels matter: they are the carries, and the carries
determine the orbit. This chapter studies a carry graph as a \emph{presentation} of its labels
\cite[ch.~3]{LM95}. Three questions come up in the problem families of
Chapter~\ref{ch:overview}:
\begin{enumerate}[label=(\arabic*)]
\item Is every infinite label sequence eventually periodic? This is the conclusion every
  certificate in the floor problems aims at.
\item Can this be decided on the graph itself, or only after making the labels deterministic?
\item Is the language of labels finitely presentable at all? The follower automaton of the
  interval problem is a carry graph with infinitely many vertices, and soficity asks whether a
  finite one would do.
\end{enumerate}
The answers are the determinization theorem (Theorem~\ref{thm:determinization}), the bridge
to regular languages (Theorem~\ref{thm:sofic-regular}) and the follower-set graph
(Section~\ref{sec:follower}).

\section{Two forms of periodicity}

\begin{definition}\label{def:allper}
A sequence \(\ell:\N\to A\) is \emph{eventually periodic}\index{eventually periodic} if \(\ell(m+L)=\ell(m)\) for some
\(L>0\) and all \(m\ge n_0\). A carry graph \(X\) has the \emph{label form}\index{eventually periodic!label form} \(\mathrm{AllPer}\) if
the label sequence of every infinite path is eventually periodic. It has the \emph{state form}\index{eventually periodic!state form}
if no vertex branches (Definition~\ref{def:branch}).
\end{definition}

The state form is face~2 of Theorem~\ref{thm:six-faces}. It speaks about vertices, while the
label form speaks about labels only.

\begin{proposition}\label{prop:state-label}
Let \(V\) be finite and \(A\) arbitrary.
\begin{enumerate}[label=(\alph*)]
\item The state form holds iff every infinite path is eventually periodic, vertices and labels
  together. In particular it implies the label form.
\item If \(X\) is right-resolving, the two forms are equivalent.
\item A homomorphism \(f:X\to Y\) carries the label form of \(Y\) back to \(X\).
\end{enumerate}
\end{proposition}

\begin{proof}
(a) Without branching, the number of vertices reachable from \(v\) is a labelled rank. This
uses only that each reach set is finite, so \(A\) need not be finite. By
Theorem~\ref{thm:evper} the path is eventually periodic. The converse is the horseshoe of
Theorem~\ref{thm:horseshoe}. (b) In a right-resolving graph the start vertex and the labels
determine the path. If the labels have period \(L\) from \(n_0\) on, the vertices
\(s(n_0+kL)\), \(k\ge0\), repeat by pigeonhole, and the repetition propagates along the
periodic labels. So the path is eventually periodic, and (a) applies. (c) \(f\) maps infinite
paths to infinite paths with the same labels.
\end{proof}

Without right-resolution the label form is strictly weaker.

\begin{example}[the \(5/3\) tube]\label{ex:tube}
At \(p/q=5/3\), shift \(\nu=\tfrac12\), primes \(\{2,3\}\), \(K=2\) and cell resolution \(k=3\),
the joint residue-and-cell graph of the interval problem has an eight-vertex strong component
with arcs
\[
0,1\xrightarrow{-2}4,5\xrightarrow{-2}6,7\xrightarrow{\;2\;}2,3\xrightarrow{\;2\;}0,1
\]
(12 arcs, data \texttt{CarryGraph/data/tube53\_k3.txt}). Vertex \(0\) has the two arcs
\(0\xrightarrow{-2}4\) and \(0\xrightarrow{-2}5\), and both targets return to \(0\). So the state
form fails. Every arc raises a level modulo 4, and its label depends only on the level, so
every infinite path reads \((-2,-2,2,2)^\infty\) from some point on, and the label form holds.
The same phenomenon occurs in a three-vertex graph, \(0\xrightarrow{a}1\), \(0\xrightarrow{a}2\),
\(1\xrightarrow{b}0\), \(2\xrightarrow{b}0\), which presents \((ab)^\infty\) and branches at
\(0\). Both checks for this graph are done by \lean{decide}.
\end{example}

This is why, in the interval problem, the state form is not monotone in the cell resolution
\(k\): refining the cells can produce two parallel tubes with the same labels.

\section{The subset construction}

\begin{definition}[{\cite[Thm.~3.3.2]{LM95}}]\label{def:det}
For \(S\subseteq V\) and \(a\in A\) let \(S\cdot a=\{w : v\xrightarrow{a}w,\ v\in S\}\). The
\emph{subset construction}\index{subset construction} \(\det X\) is the carry graph on the finite subsets of \(V\), with an
arc \(S\xrightarrow{a}S\cdot a\) whenever \(S\cdot a\neq\emptyset\).
\end{definition}

\(\det X\) is right-resolving by construction. It is computable: in Lean it is a carry graph on
\(\mathrm{Finset}\,V\). Reading a word \(w\) from \(S\) in \(\det X\) leads to the set of
vertices that \(w\) reaches from \(S\) in \(X\). Hence the language of \(\det X\) from \(\{v\}\)
is the language of \(X\) from \(v\).

\begin{lemma}[Kőnig for the subset construction]\label{lem:det-konig}
Every infinite path of \(\det X\) from a non-empty \(S_0\) lifts to an infinite path of \(X\)
from a vertex of \(S_0\), with the same labels. Neither \(V\) nor \(A\) needs to be finite.
\end{lemma}

\begin{proof}
Let \(S_0\xrightarrow{\ell_0}S_1\xrightarrow{\ell_1}\cdots\). Every vertex of \(S_m\) is reached
from \(S_0\) by a walk reading \(\ell_0\cdots\ell_{m-1}\), so for each \(m\) some vertex of the
finite set \(S_0\) starts such a walk. By pigeonhole one vertex \(v_0\) does so for every \(m\).
Among its finitely many \(\ell_0\)-successors, again one has walks reading
\(\ell_1\cdots\ell_{m}\) for every \(m\). Dependent choice gives the path.
\end{proof}

\begin{corollary}\label{cor:konig-labels}
Let \(V\) be finite and \(\ell_0\ell_1\cdots\) a label sequence every prefix of which is read by
some walk of \(X\). Then \(\ell\) is the label sequence of an infinite path of \(X\).
\end{corollary}

\begin{proof}
Run the subset construction from \(S_0=V\) along \(\ell\). Each \(S_m\) contains the end of a walk
reading \(\ell_0\cdots\ell_{m-1}\), so it is non-empty, and the \(S_m\) form a path of
\(\det X\). Lemma~\ref{lem:det-konig} lifts it.
\end{proof}

This is the step that turns a supply of finite admissible words into an obstruction: if the
prefixes of an aperiodic sequence are all read, the label form fails.

\begin{theorem}[determinization]\label{thm:determinization}\index{determinization}
For finite \(V\) the following are equivalent:
\begin{enumerate}[label=(\roman*)]
\item \(X\) has the label form;
\item \(\det X\) has the label form;
\item \(\det X\) has the state form;
\end{enumerate}
and, for finite \(A\),
\begin{enumerate}[label=(\roman*),resume]
\item \(\det X\) has polynomial growth;
\item \(h(\det X)\le 0\);
\item the label language grows polynomially: for some \(C,d\) and all \(v,n\),
  \[\#\{w : |w|=n,\ w\text{ is read from }v\}\le C(n+1)^d.\]
\end{enumerate}
\end{theorem}

\begin{proof}
(i)\(\Leftrightarrow\)(ii): a path of \(X\) gives a path of \(\det X\) with the same labels,
and Lemma~\ref{lem:det-konig} gives the converse. (ii)\(\Leftrightarrow\)(iii) is
Proposition~\ref{prop:state-label}(b), since \(\det X\) is right-resolving.
(iii)\(\Leftrightarrow\)(iv)\(\Leftrightarrow\)(v) is Theorem~\ref{thm:six-faces} for
\(\det X\). For (vi): in a right-resolving graph paths and words correspond one to one
(Proposition~\ref{prop:count}), so \(\mathrm{npaths}_{\det X}(\{v\},n)\) is the number of words
of length \(n\) from \(v\). Conversely, the words from \(S\) are the union of the words from the
vertices of \(S\). This gives \(\mathrm{npaths}_{\det X}(S,n)\le|V|\,C(n+1)^d+1\).
\end{proof}

Item (vi) is the precise meaning of the informal criterion ``\(\lambda_{\det}=1\)'' used for the
interval problem: the label language grows polynomially. It needs no Perron--Frobenius theory.
For the tube of Example~\ref{ex:tube}, (i) holds while the state form on \(X\) fails, so (iii)
holds while it fails on \(X\) itself.

\section{Languages and regularity}

\begin{definition}\label{def:sofic}
The \emph{language}\index{language} of \(X\) from \(v\) is the set \(\mathcal L(X,v)\) of words read on finite
walks from \(v\). A language \(L\subseteq A^*\) is \emph{prefix-closed}\index{prefix-closed} if it contains the empty
word and every prefix of each of its words. It is \emph{sofic}\index{sofic} if \(L=\mathcal L(X,v)\) for a
finite carry graph \(X\) and some \(v\).
\end{definition}

This is the one-sided, rooted notion. Two-sided sofic shifts appear in
Section~\ref{sec:shift}. Every \(\mathcal L(X,v)\) is prefix-closed. Reading \(X\) as a
nondeterministic automaton with start state \(v\) and every state accepting gives an NFA with
language \(\mathcal L(X,v)\). Mathlib's subset construction \lean{NFA.toDFA} run from a finite
set \(S\) on a word \(w\) ends in exactly the state that \(\det X\) reaches.

\begin{theorem}\label{thm:sofic-regular}\index{Myhill--Nerode theorem}
A language is sofic iff it is regular and prefix-closed.
\end{theorem}

\begin{proof}
A sofic language is accepted by the NFA just described, hence by a finite DFA. Conversely,
restrict a finite DFA for \(L\) to its accepting states, and keep an arc \(s\xrightarrow{a}t\)
only if \(t\) accepts. A word \(w\) is read from the start state iff all its prefixes are
accepted, and for prefix-closed \(L\) this is \(w\in L\).
\end{proof}

For \(u\in A^*\), the \emph{follower set} of \(u\) in \(L\) is
\(F_L(u)=\{w : uw\in L\}\) \cite[Def.~3.2.4]{LM95}. In Mathlib it is the left quotient
\lean{L.leftQuotient u}. Mathlib's Myhill--Nerode theorem says that \(L\) is regular iff it has
finitely many left quotients. With Theorem~\ref{thm:sofic-regular} this gives:

\begin{corollary}[{\cite[Thm.~3.2.10]{LM95}, one-sided form}]\label{cor:sofic-follower}
A prefix-closed language is sofic iff it has finitely many follower sets.
\end{corollary}

\section{The follower-set graph}\label{sec:follower}

\begin{definition}\label{def:follower-graph}
The \emph{follower-set graph}\index{follower-set graph} \(\mathcal F_A\) is the carry graph whose vertices are all
languages over \(A\), with an arc \(L\xrightarrow{a}F_L(a)\) whenever \(a\in L\).
\end{definition}

It is one graph for all languages at once, and it is right-resolving. Membership in a language
is not decidable in general, so \(\mathcal F_A\) is a tool for proofs and not for computation.

\begin{theorem}\label{thm:follower-graph}
Let \(L\) be prefix-closed.
\begin{enumerate}[label=(\alph*)]
\item Reading \(w\) from \(L\) in \(\mathcal F_A\) succeeds iff \(w\in L\), and it ends at
  \(F_L(w)\). Hence \(\mathcal L(\mathcal F_A,L)=L\), and the vertices reachable from \(L\) are
  the follower sets of the words of \(L\).
\item \(L\) is sofic iff only finitely many vertices of \(\mathcal F_A\) are reachable from
  \(L\).
\item (Minimality.) If \(X\) is right-resolving, then \(x\mapsto\mathcal L(X,x)\) is a
  homomorphism \(X\to\mathcal F_A\). For finite \(V\), the vertices reachable from
  \(\mathcal L(X,v)\) number at most \(|V|\), and at most \(2^{|V|}\) if \(X\) is arbitrary.
\item If \(A\) is finite and \(L\) is not sofic, there is an infinite sequence
  \(\ell_0\ell_1\cdots\) all of whose prefixes lie in \(L\) and have pairwise distinct follower
  sets.
\end{enumerate}
\end{theorem}

\begin{proof}
(a) By induction on \(w\), using \(F_L(aw)=F_{F_L(a)}(w)\) and the fact that \(F_L(a)\) is
prefix-closed when \(a\in L\). (b) The follower set of a word outside \(L\) is empty, so the
reachable vertices together with \(\emptyset\) are all follower sets. Apply
Corollary~\ref{cor:sofic-follower}. (c) In a right-resolving graph \(F_{\mathcal L(X,v)}(u)\) is
the language of the unique vertex that \(u\) reaches from \(v\). So the reachable follower sets
are among the \(|V|\) languages \(\mathcal L(X,x)\). For arbitrary \(X\) apply this to
\(\det X\), which has \(2^{|V|}\) states and the same language from \(\{v\}\). (d) By (b),
infinitely many vertices are reachable from \(L\), and \(\mathcal F_A\) has finite out-degree.
Geodesic Kőnig (Theorem~\ref{thm:geodesic-konig}) gives an infinite path from \(L\) whose
vertices are pairwise distinct, and by (a) these are the follower sets of its prefixes.
\end{proof}

Part~(c) merges follower-equivalent states \cite[Lemma~3.3.8]{LM95}. With
Proposition~\ref{prop:state-label}(c) it shows that a minimized certificate table is sound for
the original graph. Part (d) is the generic half of the infinite-automaton result of the
interval problem: there the follower automaton is infinite, and an injective path of states is
found. Note that (b) concerns the follower-set graph and not an arbitrary presentation. An
infinite automaton can still present a sofic language when many of its states have the same
follower set, and for the interval problem this is in fact the case.

\section{The two-sided label shift and Rauzy graphs}\label{sec:shift}

Let \(S\subseteq V\) be essential (Definition~\ref{def:core}). The words read on walks inside
\(S\) form a factorial and extendable language. Factorial means it is closed under deleting the
first or the last letter. Extendable means every word extends by a letter on each side: an
essential set gives every vertex a successor and a predecessor inside \(S\). By the duality
between languages and subshifts \cite[§1.3]{LM95}, it is the language of exactly one subshift
of \(A^{\Z}\), the \emph{label shift} \(X_S\). For \(S=\operatorname{core}X\) this is the sofic
shift presented by \(X\) \cite[Def.~3.1.3]{LM95}. If each arc is labelled by itself, it is the
edge shift \cite[Def.~2.2.5]{LM95}.

\begin{proposition}\label{prop:label-shift}\index{label shift}
For a non-empty essential \(S\) and a finite alphabet, the words of length \(n\) of \(X_S\) are
exactly the words of length \(n\) read inside \(S\). The label sequence of every bi-infinite
path inside \(S\) lies in \(X_S\). Every vertex of a bi-infinite path lies in the core.
\end{proposition}

The converse of the second statement, that every point of \(X_{\operatorname{core}}\) labels a
bi-infinite path, also holds for finite \(V\) by compactness. It is not needed here.

The \emph{Rauzy graph} of order \(n\) of a factorial language \(L\) is the carry graph on the
words of length \(n\), with an arc \(u\xrightarrow{a}v\) whenever some word \(e\) of length
\(n+1\) starts with \(u\), ends with \(v\), and has last letter \(a\). Its underlying digraph is
the Rauzy graph used in complexity theory. The Sturmian and complexity material of the interval
problem lives on these graphs.

\section*{Notes}
The subset construction and follower sets are \cite[§§3.2--3.3]{LM95}. The determinization
theorem with its growth leg comes from the interval problem: the certificate engine checks the
state form, and the tube example shows that only the det form is a property of the labels.
\cite[Thm.~21]{Oda26} is the same theorem in a different language: for a graph whose edges
carry single letters, the phase test on every strong component holds iff every infinite walk
has eventually periodic output. With this chapter it holds iff \(\det X\) has the state form
(Theorem~\ref{thm:determinization}). The graphs of \cite{Oda26} that the corpus formalises,
namely the one-step graph \(G_K(W)\), the complete graph \(G(M,K)\) and \(G'(R,K)\), all carry
single letters. Their projection map \cite[Lemma~31]{Oda26} is a carry-graph homomorphism, so
projection soundness is Proposition~\ref{prop:state-label}(c). Only the contracted graphs of
the normal form carry longer words. They enter by subdivision in
Chapter~\ref{ch:splitting}.

\section*{Lean dictionary}
\begin{center}\small
\begin{tabular}{@{}>{\raggedright}p{0.3\textwidth}>{\raggedright\arraybackslash}p{0.64\textwidth}@{}}
\toprule
Text & Lean (\texttt{CarryGraph/Labelled.lean}, \texttt{Language.lean}, \texttt{Follower.lean},
  \texttt{Shift.lean})\\
\midrule
Definition~\ref{def:allper} & \lean{EvPer}, \lean{AllPer}, \lean{AllEvPer.allPer}\\
Proposition~\ref{prop:state-label} & \lean{lnoBranch\_isLRank\_ncard}, \lean{lnoBranch\_iff\_allEvPer},
  \lean{allPer\_of\_lnoBranch}, \lean{evPer\_path\_of\_rr}, \lean{lnoBranch\_iff\_allPer\_of\_rr},
  \lean{allPer\_of\_hom}\\
Example~\ref{ex:tube} & \lean{tube}, \lean{tube53\_separates}, \lean{allPer\_of\_graded},
  \lean{diamond\_separates} (by \lean{decide})\\
Definition~\ref{def:det} & \lean{detStep}, \lean{det}, \lean{adj\_det}, \lean{det\_rightResolving},
  \lean{isWord\_det}, \lean{isWord\_det\_singleton}, \lean{words\_det\_singleton}\\
Lemma~\ref{lem:det-konig} & \lean{det\_iPath\_of\_iPath}, \lean{exists\_iPath\_of\_det},
  \lean{exists\_forall\_of\_finset}\\
Corollary~\ref{cor:konig-labels} & \lean{exists\_iPath\_of\_forall\_reads}\\
Theorem~\ref{thm:determinization} & \lean{allPer\_det\_iff}, \lean{determinization},
  \lean{allPer\_iff\_polyGrowth\_det}, \lean{allPer\_iff\_entropy\_det},
  \lean{allPer\_iff\_wordGrowth}\\
Definition~\ref{def:sofic} & \lean{langFrom}, \lean{PrefixClosed}, \lean{langFrom\_prefixClosed},
  \lean{Sofic}, \lean{sofic\_of\_finite}, \lean{langFrom\_det\_singleton}\\
NFA reading & \lean{toNFA}, \lean{mem\_evalFrom\_toNFA}, \lean{toNFA\_accepts},
  \lean{toDFA\_evalFrom\_toNFA}\\
Theorem~\ref{thm:sofic-regular} & \lean{Sofic.isRegular}, \lean{sofic\_of\_isRegular},
  \lean{sofic\_iff\_isRegular\_and\_prefixClosed}\\
Corollary~\ref{cor:sofic-follower} & \lean{sofic\_iff} (Mathlib \lean{Language.}\allowbreak
  \lean{isRegular\_iff\_finite\_range\_leftQuotient}), \lean{mem\_leftQuotient\_langFrom}\\
Definition~\ref{def:follower-graph} & \lean{followerGraph}, \lean{adj\_followerGraph},
  \lean{followerGraph\_rightResolving}\\
Theorem~\ref{thm:follower-graph} & \lean{reads\_followerGraph}, \lean{langFrom\_followerGraph},
  \lean{reach\_followerGraph}, \lean{sofic\_iff\_finite\_reach}, \lean{toFollower},
  \lean{ncard\_reach\_le\_of\_rr}, \lean{ncard\_reach\_le}, \lean{allPer\_of\_followerGraph},
  \lean{exists\_injective\_of\_not\_sofic}\\
Proposition~\ref{prop:label-shift} & \lean{blockWords}, \lean{isLanguage\_blockWords},
  \lean{labelShift}, \lean{language\_labelShift}, \lean{mem\_labelShift\_of\_path},
  \lean{mem\_core\_of\_biPath}, \lean{mem\_labelShift\_core\_of\_path}\\
Rauzy graphs & \lean{rauzy}, \lean{adj\_rauzy}, \lean{under\_rauzy}, \lean{reach\_rauzy}\\
\cite{Oda26} & \lean{CarryBridge.toWG}, \lean{CarryBridge.ofWG}, \lean{CarryBridge.odaAllPer\_iff},
  \lean{CarryBridge.allCert\_iff\_det}, \lean{CarryBridge.allCert\_iff\_wordGrowth},
  \lean{CarryBridge.odaProjHom}, \lean{CarryBridge.allPer\_of\_proj},
  \lean{CarryBridge.allCert\_oneStepU\_iff}\\
\bottomrule
\end{tabular}
\end{center}
